\documentclass[10pt,a4paper]{amsart}
\usepackage[margin=19mm]{geometry}
\usepackage{amsmath,amssymb,mathtools}
\usepackage[scr=rsfs,cal=euler]{mathalfa}
\usepackage{xfrac}
\usepackage[unicode,hidelinks,bookmarks=false]{hyperref}
\newcommand{\ie}{i.e.\ }
\newcommand{\cf}{cf.\ }
\newcommand{\resp}{respectively\ }
\newcommand{\Subsection}{\subsection}
\providecommand{\nobreakdash}{\nobreak\hbox{-}}
\setlength{\emergencystretch}{2em}
\multlinegap0pt
\usepackage{bm}
\usepackage{bbm}
\usepackage{dsfont}
\usepackage{xspace}
\usepackage{cancel}
\usepackage{mathtools}
\usepackage{amsthm}
\makeatletter
\@ifundefined{prop}{\newtheorem{prop}{Proposition}}{}
\@ifundefined{lem}{\newtheorem{lem}[prop]{Lemma}}{}
\makeatother
\newcommand{\eps}{\varepsilon}
\newcommand{\geqs}{\geqslant }
\newcommand{\leqs}{\leqslant }
\title[Siegel zeros and small gaps between zeros of the Riemann zet]{Siegel zeros and small gaps between zeros of the Riemann zeta function}
\author{Andriy Bondarenko, Winston Heap}
\date{Source: arXiv:2608.07399v1 (CC BY 4.0)}
\begin{document}
\maketitle

\noindent\emph{Source and licence.} Siegel zeros and small gaps between zeros of the Riemann zeta function. Version arXiv:2608.07399v1 (CC BY 4.0), distributed under the Creative Commons Attribution 4.0 International licence, as recorded in the arXiv licence metadata for this exact version and shown on the abstract page https://arxiv.org/abs/2608.07399v1. Full rights notice: licenses/arxiv-2608.07399-CC-BY-4.0.txt.\par\smallskip

\noindent\textit{Editorial scope.} Counted unit: the Lem whose source label is \texttt{lem:prime-L1} in the article (source0.tex lines 1109--1126, source label \texttt{lem:prime-L1}), delivered together with the article's own complete original proof of it (source0.tex lines 1128--1296, one \texttt{proof} environment of 169 lines). No mathematics is written, repaired or re-derived: statement and proof are the article's own text line for line. Required context retained uncounted: prop:density (lines 434--444). Every definition, display and label of those lines is retained unchanged. Omitted from this package and not redistributed: 1--433, 445--1108, 1127--1127, 1297--1385. Nothing inside a retained excerpt is omitted or altered: every retained body is delivered exactly as the article's lines are written, so no omission receipt of any kind accompanies this package. Citations: the retained excerpts carry no citation command at all, so no bibliography entry of the article is transcribed and no reference is redistributed. Reference adaptation: none was needed and none was made; every reference inside the delivered unit resolves inside it, so no unresolved reference or citation is produced. Printed slips kept literal: no printed word, symbol, hypothesis or number of the counted statement or of its proof was altered. Layout and class: the article's own class is \texttt{amsart} with packages including amsmath, amssymb, color, xcolor, hyperref, amsthm, amscd, geometry, mathrsfs, dsfont, mathtools, eurosym, booktabs, float; the wrapper is the lane's pinned \texttt{amsart} editorial wrapper at 10pt on A4, which restates the theorem-like environments the retained text needs and adds the article's own macro definitions that the retained text needs (\texttt{\textbackslash eps}, \texttt{\textbackslash geqs}, \texttt{\textbackslash leqs}), with extra wrapper packages (bm, bbm, dsfont, xspace, cancel, mathtools). The delivered numbering is the unit's own. Assets: no retained span contains a figure environment, a graphics-inclusion command, a PDF-page inclusion, a table or a TikZ picture; no image file is copied into this package, the package declares no asset dependency and no image file is redistributed with it. Licence: version arXiv:2608.07399v1 of this article is distributed under the Creative Commons Attribution 4.0 International licence, as recorded in the arXiv licence metadata for this exact version and shown on the abstract page https://arxiv.org/abs/2608.07399v1. A full rights notice is delivered at licenses/arxiv-2608.07399-CC-BY-4.0.txt. Characters: every retained excerpt is pure ASCII, so the delivered engine, XeLaTeX with its default Latin Modern text font, reports no missing character.
\begin{prop}
\label{prop:density}
For $1/2\leqs\sigma<1$, $H\geqs1$, and every $\eps>0$,
\begin{equation}
 \sideset{}{^*}\sum_{\psi\bmod q}N(\sigma,H,\psi)
 \ll_{\eps}
 (qH)^{7(1-\sigma)/3+\eps}
 \label{eq:density}
\end{equation}
where the sum is restricted to primitive characters modulo $q$.
\end{prop}

\begin{lem}
\label{lem:prime-L1}
Let $V$ be a
smooth function supported in $[1,2]$ and suppose that
$\|V^{(j)}\|_\infty\ll_j q^{\eps}$.  If
\begin{equation*}
 K\geqs q^{7/3+2\eta},
\end{equation*}
then
\begin{equation*}
 \sum_{\substack{\psi\bmod q\\\psi\notin\{\psi_0,\chi\}}}
 \left|\sum_n\Lambda(n)\psi(n)V(n/K)\right|
 \ll_{\eta,\eps}Kq^\eps
\end{equation*}
where $\psi_0$ is the principal character. 

% In particular, the estimate is uniform for the mild oscillatory weights produced by Lemma~\ref{lem:separation}, after its Fourier truncation exponent is chosen sufficiently small.
\end{lem}

\begin{proof}
We first replace the sum by that involving the (primitive) inducing characters, so as to apply zero-density results later. Let $\psi^*$ be the primitive character inducing $\psi$, and let
$d_\psi$ be its conductor.  If $\psi\ne\psi_0$, then $\psi^*$ is nonprincipal also. 
The difference between the two sums is 
\begin{align*}
 &\sum_n\Lambda(n)\psi(n)V(n/K)
 -\sum_n\Lambda(n)\psi^*(n)V(n/K)\\
 &\qquad=-\sum_{\substack{p\mid q\\p\nmid d_\psi}}
 \sum_{\substack{r\geqs1\\K\leqs p^r\leqs 2K}}
 (\log p)\psi^*(p^r)V(p^r/K).
\end{align*}
For a fixed prime $p$, at most two powers $p^r$ lie in $[K,2K]$ and consequently
\[
 \left|\sum_n\Lambda(n)
       (\psi(n)-\psi^*(n))V(n/K)\right|
 \ll \|V\|_\infty\sum_{p\mid q}\log p
 \ll q^{\eps}.
\]
Summing over the $\psi$ modulo $q$, the total difference is
\begin{equation*}
 O\bigl(q^{1+\eps}\bigr)=o(K).
\end{equation*}
Thus, it suffices to consider the sum
\[
 \sum_{\substack{\psi\bmod q\\\psi\notin\{\psi_0,\chi\}}}
\bigg|\sum_n\Lambda(n)\psi^*(n)V(n/K)\bigg|.
\]


We now compute the inner sum via the explicit formula. Write
\[
 \widetilde V(s)=\int_0^\infty V(x)x^{s-1}\,dx,
 \qquad
 \mathcal V_j=1+\max_{0\leqs \ell\leqs j}
                    \|V^{(\ell)}\|_\infty.
\]
From the compact support of $V$, integration by parts gives the rapid decay bounds
\begin{equation}
 |\widetilde V(\sigma+it)|
 \ll_j \mathcal V_j(1+|t|)^{-j}
 \qquad (-1\leqs\sigma\leqs2).
 \label{eq:Mellin-decay}
\end{equation}
Now, by Mellin inversion we have 
\begin{equation}
 \sum_n\Lambda(n)\psi^*(n)V(n/K)
 =\frac{1}{2\pi i}\int_{(2)}
   -\frac{L'}{L}(s,\psi^*)\widetilde V(s)\,K^sds.
 \label{eq:prime-Mellin-inversion}
\end{equation}
Shifting the contour to $\Re s=-1/2$ we encounter poles at the nontrivial zeros,
and possibly trivial zeros including the zero
at $s=0$ that occurs for an even nonprincipal character.  By
\eqref{eq:Mellin-decay}, the trivial-zero residues and the new vertical
integral contribute
\[
 O_B\bigl(\mathcal V_{B+2}\log(2qK)\bigr).
\]
Thus,
\begin{equation}
 \frac1K\left|\sum_n\Lambda(n)\psi^*(n)V(n/K)\right|
 \ll_B q^\eps
 \bigg(\sum_{\rho_{\psi^*}}
 K^{\beta-1}(1+|\gamma|)^{-B}
 +\frac{\log(2qK)}K\bigg).
 \label{eq:prime-explicit}
\end{equation}

We first estimate the tails of the zero sums and for this we need an estimate for the number of zeros. If $F$ is any
nonnegative function of a primitive character, then
\begin{equation*}
 \sum_{\psi\bmod q}F(\psi^*)
 =\sum_{d\mid q}\ \sideset{}{^*}\sum_{\xi\bmod d}F(\xi)
\end{equation*}
where the inner sum is over primitive characters modulo $d$. The
standard individual zero count
\[
 N_\xi(Y):=\#\{\rho_\xi:|\Im\rho_\xi|\leqs Y\}
 \ll Y\log(2dY) \qquad (Y\geqs1)
\]
therefore yields
\begin{equation}
 \mathcal N_q(Y):=
 \sum_{\psi\bmod q}
 \#\{\rho_{\psi^*}:|\gamma|\leqs Y\}
 \ll qY\log(2qY).
 \label{eq:aggregate-zero-count}
\end{equation}
%Both the explicit formula above and this zero count may also be found in
%\cite[Chapters~5 and~10]{IwaniecKowalski}.


Now, consider the zeros with $|\gamma|>H$ where  
\[
H=q^b,\qquad b>0
\] 
with $b$ to be determined.
Using
\eqref{eq:aggregate-zero-count} and taking $B$ large enough in terms of (any given positive) $b$ we obtain
\begin{align*}
 &\sum_{\psi\bmod q}\sum_{|\gamma|>H}
 K^{\beta-1}(1+|\gamma|)^{-B}
 \ll qH^{1-B}\log(2qH)
 \ll q^{-2}\log(2qH)
\end{align*}
since $\beta<1$. This is $o(1)$, with a fixed power saving.

For the zeros with $\beta<1/2$ we have, from
\eqref{eq:aggregate-zero-count} and the assumption $K\geqs q^{7/3+2\eta}$,
\begin{equation*}\begin{split}
 \sum_{\psi\bmod q}
 \sum_{\substack{|\gamma|\leqs H\\\beta<1/2}}K^{\beta-1}
 &\leqs K^{-1/2}\mathcal N_q(H) \\
 &\ll qHK^{-1/2}\log(2qH) \\
 &\ll q^{-1/6+b-\eta}\log(2qH)=o(1)
 \label{eq:left-zero-range}
\end{split}
\end{equation*}
provided that $b\leqs1/6$.

It remains to consider $\beta\geqs1/2$ and $|\gamma|\leqs H$.  Set
\[
 A(\sigma)=\sum_{\psi\bmod q}N(\sigma,H,\psi^*).
\]
The elementary identity
\[
 K^{\beta-1}=K^{-1/2}
 +\log K\int_{1/2}^{\beta}K^{\sigma-1}\,d\sigma
 \qquad (\beta\geqs1/2)
\]
gives, after summing one zero at a time,
\begin{align}
{\sum_{\psi\bmod q}}
   {\sum_{\substack{\rho_{\psi^*}\\
          \beta\geqs1/2,\ |\gamma|\leqs H}}}K^{\beta-1}
 &=K^{-1/2}A(1/2)
 +\log K\int_{1/2}^{1}K^{\sigma-1}A(\sigma)\,d\sigma.
 \label{eq:layer-cake-identity}
\end{align}
Now, Proposition~\ref{prop:density} gives 
\begin{equation}
 A(\sigma)
 =
 \sum_{d|q}\sideset{}{^*}\sum_{\xi\bmod d}N(\sigma,H,\xi)
 \ll (qH)^{7(1-\sigma)/3+\eps}
 \label{eq:density-in-prime-proof}
\end{equation}
on applying the divisor bound $d(q)\ll q^\eps$. 

The first term on the right of \eqref{eq:layer-cake-identity} is thus 
\[
\ll (q^{1+b})^{7/6+\eps}/q^{7/6+\eta}
\]
which is $o(1)$ provided $b<6\eta/7$, say. Using $K\geqs q^{7/3+2\eta}$ again and $H=q^b$ the integral is 
\[
\ll q^{\eps}\int_{1/2}^1 q^{-(\tfrac73+2\eta)(1-\sigma)+\tfrac73(1+b)(1-\sigma)}d\sigma
\]
which is $\ll q^\eps$, again provided $b<6\eta/7$. Thus, in total we find that 
\begin{align}
 &\frac1K
 \sum_{\substack{\psi\bmod q\\\psi\notin\{\psi_0,\chi\}}}
 \left|\sum_n\Lambda(n)\psi^*(n)V(n/K)\right|
\ll
 q^{\eps}
 +q^{1+\eps }\frac{\log(2qK)}K\ll q^\eps.
 \label{eq:primitive-prime-aggregate}
\end{align}
The result then follows on taking $0<b<\min(\tfrac16,\tfrac{6\eta}{7})$.
\end{proof}

\end{document}
